% Options for packages loaded elsewhere
% Options for packages loaded elsewhere
\PassOptionsToPackage{unicode}{hyperref}
\PassOptionsToPackage{hyphens}{url}
\PassOptionsToPackage{dvipsnames,svgnames,x11names}{xcolor}
%
\documentclass[
]{hdsr}
\usepackage{xcolor}
\usepackage{amsmath,amssymb}
\setcounter{secnumdepth}{5}
\usepackage{iftex}
\ifPDFTeX
  \usepackage[T1]{fontenc}
  \usepackage[utf8]{inputenc}
  \usepackage{textcomp} % provide euro and other symbols
\else % if luatex or xetex
  \usepackage{unicode-math} % this also loads fontspec
  \defaultfontfeatures{Scale=MatchLowercase}
  \defaultfontfeatures[\rmfamily]{Ligatures=TeX,Scale=1}
\fi
\usepackage{lmodern}
\ifPDFTeX\else
  % xetex/luatex font selection
\fi
% Use upquote if available, for straight quotes in verbatim environments
\IfFileExists{upquote.sty}{\usepackage{upquote}}{}
\IfFileExists{microtype.sty}{% use microtype if available
  \usepackage[]{microtype}
  \UseMicrotypeSet[protrusion]{basicmath} % disable protrusion for tt fonts
}{}
\makeatletter
\@ifundefined{KOMAClassName}{% if non-KOMA class
  \IfFileExists{parskip.sty}{%
    \usepackage{parskip}
  }{% else
    \setlength{\parindent}{0pt}
    \setlength{\parskip}{6pt plus 2pt minus 1pt}}
}{% if KOMA class
  \KOMAoptions{parskip=half}}
\makeatother
% Make \paragraph and \subparagraph free-standing
\makeatletter
\ifx\paragraph\undefined\else
  \let\oldparagraph\paragraph
  \renewcommand{\paragraph}{
    \@ifstar
      \xxxParagraphStar
      \xxxParagraphNoStar
  }
  \newcommand{\xxxParagraphStar}[1]{\oldparagraph*{#1}\mbox{}}
  \newcommand{\xxxParagraphNoStar}[1]{\oldparagraph{#1}\mbox{}}
\fi
\ifx\subparagraph\undefined\else
  \let\oldsubparagraph\subparagraph
  \renewcommand{\subparagraph}{
    \@ifstar
      \xxxSubParagraphStar
      \xxxSubParagraphNoStar
  }
  \newcommand{\xxxSubParagraphStar}[1]{\oldsubparagraph*{#1}\mbox{}}
  \newcommand{\xxxSubParagraphNoStar}[1]{\oldsubparagraph{#1}\mbox{}}
\fi
\makeatother


\usepackage{longtable,booktabs,array}
\newcounter{none} % for unnumbered tables
\usepackage{calc} % for calculating minipage widths
% Correct order of tables after \paragraph or \subparagraph
\usepackage{etoolbox}
\makeatletter
\patchcmd\longtable{\par}{\if@noskipsec\mbox{}\fi\par}{}{}
\makeatother
% Allow footnotes in longtable head/foot
\IfFileExists{footnotehyper.sty}{\usepackage{footnotehyper}}{\usepackage{footnote}}
\makesavenoteenv{longtable}
\usepackage{graphicx}
\makeatletter
\newsavebox\pandoc@box
\newcommand*\pandocbounded[1]{% scales image to fit in text height/width
  \sbox\pandoc@box{#1}%
  \Gscale@div\@tempa{\textheight}{\dimexpr\ht\pandoc@box+\dp\pandoc@box\relax}%
  \Gscale@div\@tempb{\linewidth}{\wd\pandoc@box}%
  \ifdim\@tempb\p@<\@tempa\p@\let\@tempa\@tempb\fi% select the smaller of both
  \ifdim\@tempa\p@<\p@\scalebox{\@tempa}{\usebox\pandoc@box}%
  \else\usebox{\pandoc@box}%
  \fi%
}
% Set default figure placement to htbp
\def\fps@figure{htbp}
\makeatother


% definitions for citeproc citations
\NewDocumentCommand\citeproctext{}{}
\NewDocumentCommand\citeproc{mm}{%
  \begingroup\def\citeproctext{#2}\cite{#1}\endgroup}
\makeatletter
 % allow citations to break across lines
 \let\@cite@ofmt\@firstofone
 % avoid brackets around text for \cite:
 \def\@biblabel#1{}
 \def\@cite#1#2{{#1\if@tempswa , #2\fi}}
\makeatother
\newlength{\cslhangindent}
\setlength{\cslhangindent}{1.5em}
\newlength{\csllabelwidth}
\setlength{\csllabelwidth}{3em}
\newenvironment{CSLReferences}[2] % #1 hanging-indent, #2 entry-spacing
 {\begin{list}{}{%
  \setlength{\itemindent}{0pt}
  \setlength{\leftmargin}{0pt}
  \setlength{\parsep}{0pt}
  % turn on hanging indent if param 1 is 1
  \ifodd #1
   \setlength{\leftmargin}{\cslhangindent}
   \setlength{\itemindent}{-1\cslhangindent}
  \fi
  % set entry spacing
  \setlength{\itemsep}{#2\baselineskip}}}
 {\end{list}}
\usepackage{calc}
\newcommand{\CSLBlock}[1]{\hfill\break\parbox[t]{\linewidth}{\strut\ignorespaces#1\strut}}
\newcommand{\CSLLeftMargin}[1]{\parbox[t]{\csllabelwidth}{\strut#1\strut}}
\newcommand{\CSLRightInline}[1]{\parbox[t]{\linewidth - \csllabelwidth}{\strut#1\strut}}
\newcommand{\CSLIndent}[1]{\hspace{\cslhangindent}#1}



\setlength{\emergencystretch}{3em} % prevent overfull lines

\providecommand{\tightlist}{%
  \setlength{\itemsep}{0pt}\setlength{\parskip}{0pt}}



 


%Graphics should all go in the figs/ directory
\graphicspath{{figs/}}
\makeatletter
\@ifpackageloaded{caption}{}{\usepackage{caption}}
\AtBeginDocument{%
\ifdefined\contentsname
  \renewcommand*\contentsname{Table of contents}
\else
  \newcommand\contentsname{Table of contents}
\fi
\ifdefined\listfigurename
  \renewcommand*\listfigurename{List of Figures}
\else
  \newcommand\listfigurename{List of Figures}
\fi
\ifdefined\listtablename
  \renewcommand*\listtablename{List of Tables}
\else
  \newcommand\listtablename{List of Tables}
\fi
\ifdefined\figurename
  \renewcommand*\figurename{Figure}
\else
  \newcommand\figurename{Figure}
\fi
\ifdefined\tablename
  \renewcommand*\tablename{Table}
\else
  \newcommand\tablename{Table}
\fi
}
\@ifpackageloaded{float}{}{\usepackage{float}}
\floatstyle{ruled}
\@ifundefined{c@chapter}{\newfloat{codelisting}{h}{lop}}{\newfloat{codelisting}{h}{lop}[chapter]}
\floatname{codelisting}{Listing}
\newcommand*\listoflistings{\listof{codelisting}{List of Listings}}
\makeatother
\makeatletter
\makeatother
\makeatletter
\@ifpackageloaded{caption}{}{\usepackage{caption}}
\@ifpackageloaded{subcaption}{}{\usepackage{subcaption}}
\makeatother
\usepackage{bookmark}
\IfFileExists{xurl.sty}{\usepackage{xurl}}{} % add URL line breaks if available
\urlstyle{same}
% fallback for those not using the hyperref driver hyperxmp:
\makeatletter
\@ifundefined{xmpquote}{\newcommand{\xmpquote}[1]{#1}}{}
\makeatother
\hypersetup{
  pdftitle={Generative AI hasn't changed learning: It's still hard to become an expert},
  pdfauthor={Ming-Wen An; Benjamin S. Baumer; Logan Dancey; Richard De Veaux; Daryl DeFord; Erika Franklin Fowler; Stephen Freund; Sorelle Friedler; Maryam Gooyabadi; Douglas Hall; Brianna Heggeseth; Nicholas J. Horton; Emmanuel Kaparakis; Antonio Laverghetta Jr.; Victoria Manfredi; Valerie Nazzaro; Thao Nguyen; Cassandra Pattanayak; Anna Plantinga; Alex Reinhart; Lynne Steuerle Schofield; Marc Schulz; Sarah Supp; Anjali Thapar; Elizabeth Upton; Zhe Wang},
  pdfkeywords={\xmpquote{Generative AI, higher education, learning
science, data science education}},
  colorlinks=true,
  linkcolor={blue},
  filecolor={Maroon},
  citecolor={Blue},
  urlcolor={Blue},
  pdfcreator={LaTeX via pandoc}}


\title{Generative AI hasn't changed learning: It's still hard to become
an expert}
\author{Ming-Wen An \and Benjamin S. Baumer \and Logan
Dancey \and Richard De Veaux \and Daryl DeFord \and Erika Franklin
Fowler \and Stephen Freund \and Sorelle Friedler \and Maryam
Gooyabadi \and Douglas Hall \and Brianna Heggeseth \and Nicholas J.
Horton \and Emmanuel Kaparakis \and Antonio Laverghetta
Jr. \and Victoria Manfredi \and Valerie Nazzaro \and Thao
Nguyen \and Cassandra Pattanayak \and Anna Plantinga \and Alex
Reinhart \and Lynne Steuerle Schofield \and Marc Schulz \and Sarah
Supp \and Anjali Thapar \and Elizabeth Upton \and Zhe Wang}
\date{}
\begin{document}

% \graphicspath{{figs/}}

\newgeometry{bottom=1.5in}


\volumeheader%
{0}%
{0}%
{00.000}

\maketitle

% Start page numbering on second page. Must appear *after* \maketitle
\thispagestyle{empty}

 % % Authors and Affiliations
 %  \begin{tabular}{cc}
 %    First Author\upstairs{\affilone,*}, Second Author\upstairs{\affilone}, Third Author\upstairs{\affilthree}
 %   \\[0.25ex]
 %   {\small \upstairs{\affilone} Affiliation One} \\
 %   {\small \upstairs{\affiltwo} Affiliation Two} \\
 %   {\small \upstairs{\affilthree} Affiliation Three} \\
 %  \end{tabular}
  
 % % Replace with corresponding author email address
 % \emails{
 %   \upstairs{*}correspondingauthor@example.edu 
 %   }
 % \vspace*{0.4in}

\begin{abstract}
Generative AI is having a substantial impact on higher education. Amid a
growing body of research about that impact, we articulate an emerging
consensus about how generative AI affects student learning that is
supported by established educational theories. Specifically, we find
that while generative AI appears to be helpful to experts, it is not
similarly helpful to novices in developing their expertise. Students
(novices) need a foundation of disciplinary knowledge to become experts
and even critical thinkers, and their use of generative AI to provide
answers, rather than help them refine their understanding, hinders their
progress on the path from novice to expert. We connect these dots,
communicate this view to various constituencies, and discuss paths
forward for higher education. While we write from our perspective as
statisticians, data scientists, and computer scientists at liberal arts
colleges, we believe that our findings will resonate across other
disciplines and institutions of higher education.
\end{abstract}

\vspace*{0.15in}

\hspace{10pt}
  \small	
  \textbf{\textit{Keywords: }} Generative AI, higher education, learning
science, data science education

\copyrightnotice

% Media summary should come next as an unnumbered section

\section{Introduction}\label{introduction}

We write as a group of 26 statistics, data science, and computer science
educators at 11 liberal arts colleges and one R1 university. We convened
in July 2026 (see Section~\ref{sec-workshop}) to discuss the
opportunities and challenges facing undergraduate higher education
brought on by the rapid adoption of generative artificial intelligence
(GenAI, hereafter) tools, and their impact upon our students. Our goal
in this paper is to present a research-informed view of how GenAI is
affecting our educational objectives -- for better and for worse -- and
to communicate that perspective succinctly to important stakeholders,
including students, colleagues, trustees, and parents. As statistics,
data science, and computer science educators from liberal arts colleges,
we focus our attention on the environments that we know best, but we
believe that what we write will resonate across much of higher
education. While we recognize that our disciplines and our institutions
do not fully represent higher education at large, we hope our careful
thinking will have value to everyone connected to these enterprises. As
liberal arts college educators, we feel well-positioned to lead on this
notable challenge, and suggest that liberal arts colleges, where faculty
are deeply engaged with both research and teaching, can serve as ideal
incubators for novel pedagogical strategies to address GenAI.

While we are critical of the most optimistic claims about GenAI's role
in undergraduate education, we are not technophobes. Neither do we have
our heads buried in the proverbial sand, nor are we afraid of change.
Most of us have doctorates in statistics, computer science, or
closely-related fields, some of us have research interests that include
subdomains of artificial intelligence, and all of us use, create, or
examine GenAI to some extent in our work. Many of us have been
recognized as classroom innovators within our disciplines. We are not of
one mind: our approaches to dealing with GenAI in our classes run the
gamut from ``let's ban it'' to ``let's embrace it.'' And yet, after
three and a half days of discussion, we arrived at a largely shared
conception of \emph{how} GenAI is currently affecting our students, our
classes, and our profession.

In what follows, we articulate that shared conception -- which is
supported by peer-reviewed research and rooted in longstanding theories
of learning -- and reflect on pedagogical alignment in the presence of
GenAI. Our analysis hinges on the differential utility of GenAI for
novices versus experts:

\begin{quote}
\textbf{Generative AI amplifies existing expertise, but its unchecked
adoption by students hinders their ability to develop that expertise.}
\end{quote}

\subsection{Our vision of a liberal arts
education}\label{our-vision-of-a-liberal-arts-education}

An undergraduate education is an investment in the student to whom it
pays lifelong dividends. Our students are the product of our classes,
not only as future professionals or scholars, but as human members of
society. We assert that the lasting artifact(s) of a collegiate course
is not the homework assignment, group project, or final exam completed
by the student, but rather the multifaceted development that occurred
within the student during the creation and process of doing the course
work. Within our respective disciplines, our responsibility as educators
is to help our students move along a path from novice to expert.

This focus on the holistic development of the student is explicit in the
mission statements of our respective institutions, which exhort us to
guide students to ``lead lives of meaning and purpose''
(\href{https://www.williams.edu/about/mission-values/}{Williams
College}), to ``educate women of promise for lives of distinction and
purpose''
(\href{https://www.smith.edu/discover-smith/history-traditions}{Smith
College}), and to ``educate our students to become autonomous thinkers,
discerning moral agents and active citizens of a democratic society''
(\href{https://catalog.denison.edu/catalog/history-mission-values/mission/}{Denison
University}), while emphasizing ``holistic individual and community
development, student leadership, individual responsibility, advocacy of
academic inquiry, and freedom of thought, opinion and expression in the
spirit of mutual respect''
(\href{https://www.wesleyan.edu/reslife/office/Mission\%20Statement1.html}{Wesleyan
University}). (All of our institutions have similar themes in their
mission statements.) While graduates of liberal arts colleges have had
excellent employment prospects relative to their peers (Gray et al.
2025), their employability is as much a by-product of a well-rounded,
humanistic education as it is an end in itself. Setting students up to
have better options and make better choices is a foreground priority.

Moreover, the skills that make liberal arts graduates such valuable
employees are not limited to the course content they absorbed. Instead,
it is their ability and determination to solve difficult problems, to
deploy multiple analytical approaches, to pose trenchant questions, to
summon the collective wisdom of previous scholars, to apply moral and
ethical reasoning from an informed perspective, to collaborate and
communicate effectively across multiple audiences with multiple
viewpoints, to adapt to new situations, and perhaps above all else, to
continue to learn on their own and to value the expertise of themselves
and others that make them such valuable employees (Holmes-Sullivan
2026). Rather than being confined to a single application, their skills
translate across adjacent domains.

The world continues to change ever more rapidly (higher education is no
different!) and at this moment in time, GenAI is in the eye of the
storm. Other moments in time brought the advent of the Internet and the
arrival of cell phones, and we do not know what future developments will
require. We convened to discuss the impact of GenAI on the teaching of
statistics, data science, and computer science at liberal arts colleges
in the present moment, and left convinced that while GenAI will disrupt
some of our practices (as other inventions have), it will not alter the
fundamental relationship between our students and the world that we have
just described. Just as with other new technological changes, we will
continue to help our students develop as whole persons belonging to a
lineage and community of thinkers, and prepare them for the unknowable
challenges of tomorrow, as best we can.

To do so requires an in-depth conversation about how generative AI
impacts learning.

\section{How is generative AI affecting
learning?}\label{how-is-generative-ai-affecting-learning}

\subsection{What is learning?}\label{what-is-learning}

We know from the learning sciences that what people need to understand
their world is a working mental model of the systems with which they
interact (DiSessa 1982; Schauble et al. 1991). This is true for
everyone: the doctor and the medical assistant, the receptionist and the
executive, the plumber and the chemist. These mental models allow people
to employ causal reasoning, such as making a prediction about what will
happen and then testing that prediction against reality, an act that may
result in an update to the mental model. The more accurate the mental
model of the system, the more competent the professional is.

Moreover, humans develop expertise by committing their working
understanding (i.e., schemas) to long-term memory, such that routine
tasks become automated (Van Merriënboer and Sweller 2010). Learning is
the result of this automation. Think of the way you can brush your
teeth, tie your shoes, or work the foot pedals in a car without thinking
about what you are doing, or further, while thinking about something
else entirely. Building these automated schemas requires deliberate,
repeated practice (Ericsson and Pool 2017).

In statistics and data science, we often work with a variety of systems
that include the data context, statistical models (e.g., a linear
regression model\footnote{\href{https://en.wikipedia.org/wiki/Linear_regression}{Linear
  regression} is a flexible and powerful way to predict an outcome as a
  function of a set of predictor variables.}), and computational
software (e.g., R or Python). Because fitting and interpreting linear
regression models is a standard practice in scientific literature, GenAI
can produce many factually accurate statements about linear regression
models. It can also write the R code necessary to fit and visualize
these models. We are not blind to the fact that it can do nearly all of
the tasks that we ask our introductory statistics students to do (DeLuca
et al. 2025; {Lu et al.} 2025).

What then is the value of learning statistics when GenAI can ``do data
analysis'' for you? Let's first recognize that there is a difference
between a person who understands statistics (i.e., has a working mental
model of regression committed to long-term memory) and a person who can
use GenAI to perform statistical tasks (i.e., knows how to prompt AI to
get the requested information). We might call the former ``human
intelligence'' and the latter ``human-guided artificial intelligence''.

An analogy to music might help. GenAI amplifies existing skills (Wickham
2026) the same way an incredible new amp could have expanded Jimi
Hendrix's or B. B. King's (possibly two of the greatest guitarists to
ever play) sounds and techniques in a way that would have been
prohibitively expensive in terms of time, money, or both when they were
alive. However, for the novice who is just learning to play, having a
fancy amp might be cool, but it doesn't really help them play any
better. The evidence for the efficacy of GenAI-based tutoring is mixed
at best (Kestin et al. 2025; Bastani et al. 2025), and there is simply
no substitute for good old-fashioned practice (Ericsson and Pool 2017).
Without an automated mental model of scales, chord progressions, timing,
and without the physical ability to repeat accurate fingerings and
strumming, the novice will never be able to play in a band, let alone
improvise and jam the way Hendrix or King could. Even a promising high
school or college musician with the new amp won't be at the level of
Hendrix or King \emph{without simultaneous and rigorous dedication to
practice and learning}. Nobody gets better at playing a musical
instrument\footnote{We chose the guitar-playing analogy for alignment
  and accessibility, but considered many others, including playing the
  violin, hitting a baseball, composing a photograph,
  drawing/painting/sculpting a work of art, crafting a piece of
  furniture, scaling a mountain, etc. These are all activities in which
  expensive equipment can help experts achieve their goals more
  efficiently, but will not help novices achieve theirs.} by sitting on
the couch---you have to put in the work, and the work is hard.

We also know from the learning sciences that while ``critical thinking''
is perhaps the most durable skill that we can help our students develop,
it is really a complex collection of skills that cannot be developed in
the absence of methodological content and subject matter context
(Willingham 2008; Pithers and Soden 2000; Tricot and Sweller 2014). That
is, students need both fundamental disciplinary knowledge and domain
expertise in order to apply critical thinking. Thus, it's not possible
to think critically about a regression model for home sales until you
know something about regression models \emph{and} something about the
real estate market. You can't be
\href{https://en.wikipedia.org/wiki/Symphony_No._94_(Haydn)}{surprised
by a symphony} until you've developed a mental model for what symphonies
are supposed to sound like. You can't be fooled by a curveball until
you've learned to hit a fastball by anticipating its trajectory and
coordinating your swing to meet that expectation.

Proponents of GenAI would have you believe that interacting with a
chatbot can help you build these foundational skills, but these claims
would require strong experimental evidence (which is scant (Gerlich
2025)), and if you accept our analogy, they defy logic. As educators, we
have been helping students develop foundational skills much in the
manner that music teachers have: first we explain, then we demonstrate,
then students practice and receive feedback. As any music teacher will
tell you, this process doesn't work if students don't build muscle
memory (i.e., automate their schema) through practice. For many
students, the use of GenAI is impeding that practice -- not enhancing it
-- by giving them a largely effortless way to accomplish tasks without
actually refining their mental models and muscle memories.

\subsection{Early returns on learning with generative
AI}\label{early-returns-on-learning-with-generative-ai}

The students who most recently graduated from our institutions began
their college experience shortly before ChatGPT kicked off the current
wave of chatbot-based GenAI. If the loftiest claims about GenAI were
true, we might expect that student work would have improved noticeably
in this period. However, we can say anecdotally, but confidently, that
we are not seeing such an improvement in what they are submitting. While
it is difficult to untangle all of the cohort effects these students
have experienced---growing up with smartphones and social media, the
COVID-19 pandemic disrupting their high school experience, and now the
rapid expansion of GenAI---lessons from the learning sciences can help
us unpack the real-world impact of GenAI on student learning processes.

A growing body of research supports the conclusion that for a variety of
reasons, while GenAI can help experts leverage their knowledge and
experience to perform tasks faster and more efficiently, the opposite is
true for novices (Daniotti et al. 2026). By definition, novices don't
have the background knowledge that prepares them to ask the right
questions or identify the structure of a problem, and thus chatbots
often lead them on a circuitous path to an answer to the wrong question.
Because they lack the knowledge that allows them to accurately judge the
veracity of GenAI output, novices are generally ill-equipped to navigate
chatbot responses that produce false statements. Novices also lack the
experience that helps experts quickly triage promising and doomed
suggested routes to solve a problem. Moreover, novices suffer from
metacognitive failures that experts escape: they are more likely to
believe that conversations with GenAI are deepening their understanding,
even while their assessed performance is declining (Prather et al. 2023,
2024; Fan et al. 2025). These insights help to explain why novices often
end up chasing their tail when using GenAI.

Furthermore, reliance on GenAI accustoms novices with habits that erode
the development of the skills of liberal arts graduates we identified
above that set them up for success post-college. GenAI's ability to
provide quick, easy answers to difficult problems is alluring to novices
in the short-term, but it may impede their ability to build persistence,
a far more important quality for their long-term success (Liu et al.
2026). Novices who rely heavily on GenAI assistance may also fail to
cultivate the mentoring and peer relationships that reinforce learning
and propel them into future career success (Silva 2025; Hou et al.
2025)---relationships which experts have already had the time and
experience to develop. As they become more dependent on GenAI, students
concomitantly become less likely to seek out human-centric supports for
learning (e.g., office hours, tutoring, study groups, etc.), leaving
them even more isolated and disconnected from their learning community
(Hou et al. 2025; Glickman et al. 2026).

This disjunction in the usefulness of GenAI between experts and novices
offers an explanation of how two seemingly contradictory observations
can simultaneously be true. GenAI \emph{is} useful to those with
experience and training, and it may very well lead to a realization of
many of the claims being made by AI companies and its proponents.
Long-stubborn
\href{https://www.nytimes.com/2026/07/23/science/fields-medals-pardon-deng-wang-tsimerman.html}{math
problems \emph{are} being solved} thanks to assistance from GenAI.
\href{https://www.nytimes.com/2026/04/02/technology/ai-billion-dollar-company-medvi.html}{New
possibilities for entrepreneurship} \emph{are} being fueled by GenAI.
And yet the critics of AI also stand on firm ground. Rampant,
ill-advised use of GenAI is \emph{not} helping a generation of students
build the very expertise that they will need to benefit from this newly
ubiquitous technology, even as we and they are well-aware of its growing
role in their future careers (Imundo et al. 2024).

Thus, GenAI is working primarily as an amplifier (or multiplier) of the
skills that one already has.

We see mounting evidence that GenAI is hindering our students' ability
to move along the path from novice to expert, thus delaying, or even
preventing, their ability to access both the level of students'
historical achievements and growth as well as the increased productivity
that GenAI promises. Students are beginning to realize this (Gaines
2026; Schwartz and Diliberti 2026), and we observe increasing signs of
their resentment towards GenAI, in the
\href{https://www.npr.org/2026/05/20/nx-s1-5822419/ai-colleges-commencement-booing}{heckling
of graduation speakers}, the formation of
\href{https://oberlinreview.org/35718/opinions/luddite-club-implores-oberlin-to-opt-out-of-ai/}{student
clubs} and
\href{https://www.bloomberg.com/news/articles/2026-05-19/ai-on-college-campuses-sparks-pushback-protests-booing-at-graduation}{protests
against AI}, and outright refusal to use AI (McMurtrie 2025). Rather
than using GenAI because they see it as the best tool for their
learning, many students report feeling compelled (Yang 2026) to use
GenAI due to shifting norms, competition with peers, stress,
desperation, or what we view as a misguided attempt to prepare
themselves for a workplace in which they will inevitably be asked to use
GenAI.

\section{What our key constituencies need to know about learning with
GenAI}\label{what-our-key-constituencies-need-to-know-about-learning-with-genai}

In this rapidly changing landscape, we hope that this document will help
all of us (students, faculty, trustees, parents, administrators, alums,
and industry professionals) re-center our perspectives on higher
education and GenAI's role within it (Klopfer and Madden 2026). This
will require realignment from all parties, most notably including
students and faculty. The remainder of this paper addresses four key
constituencies separately, building upon the state-of-the-world we have
articulated in this Introduction.

\subsection{Students}\label{sec-students}

We want our students to understand that we care about them, their
future, and their learning. We don't want to have adversarial
relationships with them as the cheating police. As articulated above,
our goal is to move them along the path from novice to expert and help
them develop as humans.

The evidence and our experience tell us that in order to develop
expertise, students need to build accurate mental models of the systems
they are studying atop a foundation of facts, and then continue to
refine and adorn those models through repeated practice and exposure to
new material. This practice has many forms, both newfangled and
old-fashioned, including homework problem sets, in-class lectures and
activities, and textbook readings. This practice may be difficult and
time consuming, but it is a necessary component of the learning process.

While GenAI may be useful for some students in some cases, especially in
more advanced courses where some expertise has already been developed,
the evidence is mounting that the vast majority of GenAI use among
college students is hindering their ability to build and refine their
understanding of the material they are studying, which in turn, perhaps
ironically, leaves them \emph{less} well-prepared to access the benefits
of GenAI that experienced professionals report (Gates 2026).

As a consequence, we are re-orienting our courses around assessments
that take place in AI-free environments: oral presentations and exams,
in-class pencil-and-paper quizzes and tests. We want our students to
understand that this shift in assessment aligns with our learning goals
for them, and is designed to best prepare them for success in their
future endeavors, which we acknowledge will likely involve GenAI.

Lastly, we want our students to come to class and office hours, work
with their classmates, and engage fully in the challenging process of
learning. We encourage them to be active partners in their education and
cultivate strong relationships with faculty and peers. The community we
build creates lifelong bonds while supporting the learning process and
unlocking unexpected future paths. Ultimately, these interpersonal
dynamics foster human-centered capabilities that GenAI simply cannot
replicate.

\subsection{Parents}\label{parents}

We want parents to understand that we know why they are asking hard
questions about the value of a college education amid the current
headwinds of skyrocketing costs, an uncertain job market, and new
technology (including GenAI) that is rapidly changing virtually every
industry. We know that parents want their children to be ``successful''
in life (however that is defined). As highlighted in our mission
statements, we affirm that our goal is to set their children up for
adult lives in which they have better choices and make better decisions.
We can do this by filling their minds, bodies, and spirits with
knowledge, the ability to think critically about the world, and the
courage to apply moral and ethical reasoning to the unknowable
challenges they will face.

We understand that most of our graduates will be hired into positions in
which they will be expected to leverage GenAI in the workplace. But
because of the argument we have laid out above, we don't believe that
allowing them unfettered access to GenAI during their undergraduate
years is the best way to prepare them for that future. To the contrary,
we believe that the best way to prepare them for a future in which they
will use GenAI productively is to concentrate our efforts on leveraging
our expertise to develop their mental models, primarily in AI-free
environments. We believe that welcoming them, as whole persons, into our
existing communities of scholarship and learning is the best way to
position them for a lifetime of success, happiness, and gratification.

Encouraging your children to be fully present, active partners in their
education (see Section~\ref{sec-students}) is the best way to help them
navigate this challenging period of their lives.

\subsection{Colleagues}\label{colleagues}

We want our colleagues to know that we share their frustration with and
worry about the challenges that GenAI is bringing to our jobs. We
acknowledge that GenAI is having a profound impact on our teaching, our
relationships with our students, and our sense of satisfaction with our
jobs. We, too, experience moments of deep unease about the contours of
the future of higher education, and what that means for us as teachers,
scholars, and employees.

However, we also want our colleagues to understand that we see a
research-driven consensus emerging about the way that GenAI is affecting
the learning that we facilitate in our courses. As we gain a clearer
understanding of the role that GenAI plays in higher education, we see
several paths forward that, in the best-case scenario, could lead to a
reinvigoration of the liberal arts mission and leave us in a place where
students and faculty are aligned on their respective goals about
learning. It is up to us to delineate the boundaries of GenAI's role in
that world and to conduct the education research that shows how it can
and cannot productively be used. If we fail to act, others will do this
work for us. While we disagree (constructively and collegially) about
how best to adapt our courses to GenAI, we \emph{do} agree about what is
going on and what we want for our students, and we hope this unity can
sustain us as we tackle the oncoming challenges. Continued active
communication among us can facilitate the novel tactics and techniques
we will need.

\subsection{Trustees and
administrators}\label{trustees-and-administrators}

We want our trustees and administrators to know that we are grateful for
their support of our institutions and our core educational missions. We
recognize the value of their experiences beyond academia and our joint
commitment to educating our students. We value what they love about our
liberal arts colleges and the education they received: deep and critical
engagement with complex and challenging ideas, the ability to express
oneself fully, lifelong friendships with fellow students and mentorship
from professors, and the cultivation of a love of learning. These core
values at the center of a liberal arts education endure even as
technology continues to evolve.

Our aim -- as we consider how best to teach our students -- is to
preserve these values and the spirit of learning that our institutions
cultivate. We're confident that future generations of graduates can and
will go on to lives of purpose and meaning, successful careers, and
important leadership positions. Because of the research-informed view
described above, we believe that in order for students to realize their
potential they must move towards mastering the core ideas of our
subjects without the compromises to their learning brought on by GenAI.

We recognize that trustees and administrators are hearing from all sides
that GenAI is revolutionizing many industries, and that they are feeling
pressure to adapt our curricula to prepare graduates to enter that
world. We \emph{are} doing this adaptive work, and we won't abdicate our
responsibility and authority to determine learning goals and curricula
to other constituencies. As teacher-scholars and career educators, we
need the support of our trustees and administrators now more than ever
as we meet the new challenges and opportunities that accompany GenAI.

\section{Conclusion and Paths
Forward}\label{conclusion-and-paths-forward}

For many students, the use of GenAI early in their educational careers
is impeding -- not enhancing -- their development from novice to expert.
GenAI gives them a largely effortless way to accomplish tasks that
diminishes their ability to refine their mental models of complex
systems. Reliance on GenAI decreases their efforts in building learning
communities (professors, peers, and friends) that will support them in
the long term. Using GenAI to quickly answer hard problems that would
have previously taken them hours, days, or even weeks to solve, stunts
the development of their persistence and endurance that they will likely
need in their careers and lives post-schooling.

In fairness, it is likely we, teachers, have failed to impress upon
students this philosophy of practice and the benefits of educational
communities. Students may feel that some of our more time-consuming
assignments (e.g., problem sets with by-hand calculations) are pointless
``busywork''. We must do more to both align our assignments with our
learning goals, and to convince our students that practice, especially
with their peers and with us, while occasionally repetitive, is a
necessary component of learning. The question then, is how we redirect
student behavior so that sufficient practice continues to occur.

We also may need to rethink how we distribute time in relation to the
assessments that we give. \emph{All of us} are increasingly interested
in implementing in-person assessments that take place in AI-free
environments (e.g., oral exams, in-class proctored exams, etc.).
Emphasizing in-classroom assessments balanced with fewer summative
assessments outside of class raises challenges about effective teaching
strategies. As we think through when to assess, we must also reimagine
what we assess. Just as our classwork and homework assignments should
align to our learning goals, so too must our assessments match that
which we claim to believe is important.

Ultimately, our imperative is to re-incentivize the student behaviors
that will help them move from novice to expert so that they can
eventually use GenAI and/or any future technological advancement to
amplify their skills as intellectually independent contributors to the
wider world.

\subsection{Acknowledgements}\label{sec-workshop}

This paper was incubated at the
\href{https://beanumber.github.io/aalac2026/}{\emph{Generative AI and
Data Analysis: implications for data science curriculum and pedagogy}
workshop} sponsored by the Alliance to Advance Liberal Arts Colleges
(AALAC) and held at Wesleyan University from June 19-22, 2026.

Workshop participants were provided with a reading list of nearly 60
items, most of which were recently published studies about GenAI in
higher education. During the workshop, participants heard from various
speakers about how GenAI use was affecting higher education. This
included a plenary talk given by Alex Reinhart of Carnegie Mellon
University entitled ``The role of AI in statistics and data science
education.'' We also heard from Nicole Stanton (Provost and Senior
Vice-President for Academic Affairs at Wesleyan) about the
administration's perspective on GenAI, Rachel Barlow (Director for the
Center for Faculty Career Development at Wesleyan) about how students
are using and feeling about GenAI, Casey Pattanayak (Jack and Sandra
Polk Guthman '65 Director of the Quantitative Analysis Institute at
Wellesley) about student use of AI in statistics and data science
courses, Zhe Wang (Denison) about AI Sentiment in the Classroom, and
Alex Reinhart again about the differences between student writing and
LLM writing. Each day, participants reacted to written statements about
GenAI, and those results were shared with the group the following day,
which helped us build consensus around some ideas while laying bare a
lack of consensus on others. Many conversations, breakout sessions, and
debates occurred in both formal and informal settings.

Participants included: Ming-Wen An, Benjamin Baumer, Logan Dancey,
Richard De Veaux, Daryl DeFord, Erika Franklin Fowler, Stephen Freund,
Sorelle Friedler, Maryam Gooyabadi, Douglas Hall, Brianna Heggeseth,
Nicholas Horton, Emmanuel Kaparakis, Antonio Laverghetta Jr., Victoria
Manfredi, Valerie Nazzaro, Thao Nguyen, Pavel Oleinikov, Cassandra
Pattanayak, Anna Plantinga, Alex Reinhart, Jennifer Rose, Lynne Steuerle
Schofield, Marc Schulz, Sarah Supp, Anjali Thapar, Elizabeth Upton, and
Zhe Wang.

We are grateful to Martha Besade for extensive administrative and
logistical support for the workshop.

\section*{References}\label{references}
\addcontentsline{toc}{section}{References}

\protect\phantomsection\label{refs}
\begin{CSLReferences}{1}{1}
\bibitem[\citeproctext]{ref-bastani2025generative}
Bastani, Hamsa, Osbert Bastani, Alp Sungu, Haosen Ge, Özge Kabakcı, and
Rei Mariman. 2025. {``Generative {AI} Without Guardrails Can Harm
Learning: Evidence from High School Mathematics.''} \emph{Proceedings of
the National Academy of Sciences} 122 (26): e2422633122.
\url{https://doi.org/10.1073/pnas.2422633122}.

\bibitem[\citeproctext]{ref-daniotti2026using}
Daniotti, Simone, Johannes Wachs, Xiangnan Feng, and Frank Neffke. 2026.
{``Who Is Using {AI} to Code? Global Diffusion and Impact of Generative
{AI}.''} \emph{Science}, 831--35.
\url{https://doi.org/10.1126/science.adz9311}.

\bibitem[\citeproctext]{ref-deluca2025developing}
DeLuca, Laura S, Alex Reinhart, Gordon Weinberg, Michael Laudenbach,
Sydney Miller, and David West Brown. 2025. {``Developing Students'
Statistical Expertise Through Writing in the Age of {AI}.''}
\emph{Journal of Statistics and Data Science Education} 33 (3): 266--78.
\url{https://doi.org/10.1080/26939169.2025.2497547}.

\bibitem[\citeproctext]{ref-disessa1982unlearning}
DiSessa, Andrea A. 1982. {``Unlearning Aristotelian Physics: A Study of
Knowledge-Based Learning.''} \emph{Cognitive Science} 6 (1): 37--75.
\url{https://doi.org/10.1207/s15516709cog0601_2}.

\bibitem[\citeproctext]{ref-ericsson2016peak}
Ericsson, Anders, and Robert Pool. 2017. \emph{Peak: Secrets from the
New Science of Expertise}. HarperOne.

\bibitem[\citeproctext]{ref-fan2025beware}
Fan, Yizhou, Luzhen Tang, Huixiao Le, et al. 2025. {``Beware of
Metacognitive Laziness: Effects of Generative Artificial Intelligence on
Learning Motivation, Processes, and Performance.''} \emph{British
Journal of Educational Technology} 56 (2): 489--530.
\url{https://doi.org/10.1111/bjet.13544}.

\bibitem[\citeproctext]{ref-gaines2026ai}
Gaines, Lee V. 2026. \emph{This Big University System Is Embracing {AI}.
Students and Faculty Aren't All on Board}. NPR.
\url{https://www.npr.org/2026/05/25/nx-s1-5772820/artificial-intelligence-education-technology-california-state-university}.

\bibitem[\citeproctext]{ref-gates2026ai}
Gates, Bill. 2026. \emph{The Turbulent {AI} Era Is Here. The Choices We
Make Now Are Critical.} GatesNotes.
\url{https://www.gatesnotes.com/work/make-ai-work-for-everyone/reader/expanding-access-to-health-care-through-ai}.

\bibitem[\citeproctext]{ref-gerlich2025ai}
Gerlich, Michael. 2025. {``AI Tools in Society: Impacts on Cognitive
Offloading and the Future of Critical Thinking.''} \emph{Societies} 15
(1): 6. https://doi.org/\url{https://doi.org/10.3390/soc15010006}.

\bibitem[\citeproctext]{ref-glickman2026tutor}
Glickman, Mark, Jun Yan, Nalisnick Eric, and Mine Çetinkaya-Rundel.
2026. {``A Tutor, a Tool, and a Sounding Board: How Statistics Students
Use Generative {AI}.''} \emph{CHANCE} 39 (2): 39--44.
\url{https://doi.org/10.1080/09332480.2026.2698398}.

\bibitem[\citeproctext]{ref-gray2025measuring}
Gray, Kyle, Daniel Rossman, Elizabeth Davidson Pisacreta, Catharine Bond
Hill, and Madeline Trimble. 2025. {``Measuring the Economic Value of a
Liberal Education.''} \emph{ITHAKA S+R}, ahead of print.
\url{https://doi.org/10.18665/sr.323160}.

\bibitem[\citeproctext]{ref-holmes2026ai}
Holmes-Sullivan, Robin. 2026. \emph{Why {AI} Makes a Liberal Arts
Education Even More Invaluable}. U.S. News and World Report.
\url{https://www.usnews.com/opinion/articles/2026-01-20/college-ai-education-career-liberal-arts-opinion}.

\bibitem[\citeproctext]{ref-hou2025all}
Hou, Irene, Owen Man, Kate Hamilton, et al. 2025. {``'All Roads Lead to
ChatGPT': How Generative {AI} Is Eroding Social Interactions and Student
Learning Communities.''} \emph{Proceedings of the 30th ACM Conference on
Innovation and Technology in Computer Science Education v. 1}, 79--85.
\url{https://doi.org/10.1145/3724363.3729024}.

\bibitem[\citeproctext]{ref-imundo2024expert}
Imundo, Megan N, Micah Watanabe, Andrew H Potter, Jiachen Gong, Tracy
Arner, and Danielle S McNamara. 2024. {``Expert Thinking with Generative
Chatbots.''} \emph{Journal of Applied Research in Memory and Cognition}
13 (4): 465. \url{https://doi.org/10.1037/mac0000199}.

\bibitem[\citeproctext]{ref-kestin2025ai}
Kestin, Greg, Kelly Miller, Anna Klales, Timothy Milbourne, and Gregorio
Ponti. 2025. {``{AI} Tutoring Outperforms in-Class Active Learning: An
{RCT} Introducing a Novel Research-Based Design in an Authentic
Educational Setting.''} \emph{Scientific Reports} 15 (1): 17458.
\url{https://doi.org/10.1038/s41598-025-97652-6}.

\bibitem[\citeproctext]{ref-mit2026ai}
Klopfer, Eric, and Sam Madden. 2026. \emph{Report of {MIT}'s Ad Hoc
Committee on AI Use in Teaching, Learning, and Research Training}.
Massachusetts Institute of Technology.
\url{https://aiandeducation.mit.edu/report/}.

\bibitem[\citeproctext]{ref-liu2026ai}
Liu, Grace, Brian Christian, Tsvetomira Dumbalska, Michiel A Bakker, and
Rachit Dubey. 2026. {``{AI} Assistance Reduces Persistence and Hurts
Independent Performance.''} \emph{arXiv Preprint arXiv:2604.04721},
ahead of print. \url{https://doi.org/10.48550/arXiv.2604.04721}.

\bibitem[\citeproctext]{ref-lu2025stateval}
{Lu, Yuchen, Run Yang, Yichen Zhang, et al.} 2025. {``{StatEval}: A
Comprehensive Benchmark for Large Language Models in Statistics.''}
\emph{arXiv Preprint arXiv:2510.09517}, ahead of print.
\url{https://doi.org/10.48550/arXiv.2510.09517}.

\bibitem[\citeproctext]{ref-mcmurtrie2025ai}
McMurtrie, Beth. 2025. \emph{Teaching: How to Respond When Students
Don't Want to Work with {AI}}. The Chronicle of Higher Education.
\url{https://www.chronicle.com/newsletter/teaching/2025-12-11}.

\bibitem[\citeproctext]{ref-pithers2000critical}
Pithers, Robert T, and Rebecca Soden. 2000. {``Critical Thinking in
Education: A Review.''} \emph{Educational Research} 42 (3): 237--49.
\url{https://doi.org/10.1080/001318800440579}.

\bibitem[\citeproctext]{ref-prather2023s}
Prather, James, Brent N Reeves, Paul Denny, et al. 2023. {``{`It's Weird
That It Knows What i Want'}: Usability and Interactions with Copilot for
Novice Programmers.''} \emph{ACM Transactions on Computer-Human
Interaction} 31 (1): 1--31. \url{https://doi.org/10.1145/3617367}.

\bibitem[\citeproctext]{ref-prather2024widening}
Prather, James, Brent N Reeves, Juho Leinonen, et al. 2024. {``The
Widening Gap: The Benefits and Harms of Generative {AI} for Novice
Programmers.''} \emph{Proceedings of the 2024 ACM Conference on
International Computing Education Research-Volume 1}, 469--86.
\url{https://doi.org/10.1145/3632620.3671116}.

\bibitem[\citeproctext]{ref-schauble1991causal}
Schauble, Leona, Robert Glaser, Kalyani Raghavan, and Miriam Reiner.
1991. {``Causal Models and Experimentation Strategies in Scientific
Reasoning.''} \emph{The Journal of the Learning Sciences} 1 (2):
201--38. \url{https://doi.org/10.1207/s15327809jls0102_3}.

\bibitem[\citeproctext]{ref-schwartz2026more}
Schwartz, H, and M Diliberti. 2026. \emph{More Students Use {AI} for
Homework, and More Believe It Harms Critical Thinking}. Rand
Corporation.
\url{https://www.rand.org/pubs/research/_reports/RRA4742-1.html}.

\bibitem[\citeproctext]{ref-silva2025ai}
Silva, Elise. 2025. \emph{University Students Feel 'Anxious, Confused
and Distrustful' about {AI} in the Classroom and Among Their Peers}. The
Conversation.
\url{https://theconversation.com/university-students-feel-anxious-confused-and-distrustful-about-ai-in-the-classroom-and-among-their-peers-258665}.

\bibitem[\citeproctext]{ref-tricot2014domain}
Tricot, André, and John Sweller. 2014. {``Domain-Specific Knowledge and
Why Teaching Generic Skills Does Not Work.''} \emph{Educational
Psychology Review} 26 (2): 265--83.
\url{https://doi.org/10.1007/s10648-013-9243-1}.

\bibitem[\citeproctext]{ref-van2010cognitive}
Van Merriënboer, Jeroen JG, and John Sweller. 2010. {``Cognitive Load
Theory in Health Professional Education: Design Principles and
Strategies.''} \emph{Medical Education} 44 (1): 85--93.
\url{https://doi.org/10.1111/j.1365-2923.2009.03498.x}.

\bibitem[\citeproctext]{ref-wickham2026ai}
Wickham, Hadley. 2026. \emph{Y Code When Ai?} Substack: \emph{Tidy
design principles}.
\url{https://tidydesign.substack.com/p/y-code-when-ai}.

\bibitem[\citeproctext]{ref-willingham2008critical}
Willingham, Daniel T. 2008. {``Critical Thinking: Why Is It so Hard to
Teach?''} \emph{Arts Education Policy Review} 109 (4): 21--32.
\url{https://doi.org/10.3200/AEPR.109.4.21-32}.

\bibitem[\citeproctext]{ref-yang2026ai}
Yang, Michelle. 2026. \emph{{AI}'s Impact on Education at CMU and
Beyond}. The Tartan.
\url{https://the-tartan.org/2026/02/09/ais-impact-on-education-at-cmu-and-beyond/}.

\end{CSLReferences}




\end{document}
